% BHU PYQ -- Transcribed & Corrected
% Paper: GGRSE21 / GGRSE21: Basics of Remote Sensing
% Exam: B.A. / B.Sc. Semester II Examination, 2024-25 (NEP)
% Category: Skill Enhancement Course (SEC)
% Department: Geography

\documentclass[11pt,a4paper]{article}
\usepackage[a4paper,top=2.5cm,bottom=2.5cm,left=2.8cm,right=2.8cm,headheight=14pt]{geometry}
\usepackage{amsmath,amssymb}
\usepackage[T1]{fontenc}
\usepackage[utf8]{inputenc}
\usepackage{lmodern,microtype,fancyhdr,enumitem,hyperref,lastpage,parskip}

\hypersetup{
    colorlinks=true,
    urlcolor=black,
    linkcolor=black,
    pdftitle={GGRSE21: Basics of Remote Sensing -- BHU Sem II 2024-25 (NEP)}
}

\pagestyle{fancy}
\fancyhf{}
\renewcommand{\headrulewidth}{0.4pt}
\renewcommand{\footrulewidth}{0.4pt}
\fancyhead[L]{\small\textit{Banaras Hindu University}}
\fancyhead[R]{\small\textit{GGRSE21: Basics of Remote Sensing (Sem II 2024-25)}}
\fancyfoot[C]{\small Page \thepage\ of \pageref{LastPage}}

\newlist{parts}{enumerate}{2}
\setlist[parts,1]{label=(\alph*),leftmargin=*,itemsep=4pt,topsep=2pt}
\setlist[parts,2]{label=(\roman*),leftmargin=*,itemsep=2pt,topsep=2pt}
\newcommand{\pts}[1]{\hfill\textit{[#1]}}

\begin{document}

\begin{center}
  {\large\bfseries BANARAS HINDU UNIVERSITY}\\[2pt]
  {\normalsize B.A. / B.Sc. Semester II Examination, 2024-25 (NEP)}\\[6pt]
  \rule{0.8\linewidth}{0.5pt}\\[6pt]
  {\large\bfseries GEOGRAPHY}\\[4pt]
  {\large\bfseries Paper Code: GGRSE21 : Basics of Remote Sensing}\\[4pt]
  {\normalsize Time: 2:30 Hours\hfill Maximum Marks: 70}
\end{center}

\rule{\linewidth}{0.4pt}

\smallskip
\noindent\textit{Instructions: Attempt \textbf{five} questions including Question 1 which is compulsory. Answer each part of Question 1 in approximately 200 words and of remaining questions in approximately 1000 words. The figures in the right-hand margin indicate marks. Illustrate your answers with suitable maps and diagrams.}\\[4pt]
\noindent\textit{निर्देश: प्रश्न 1 सहित पाँच प्रश्न हल कीजिए, जो अनिवार्य है। प्रश्न 1 के प्रत्येक भाग का उत्तर लगभग 200 शब्दों में और शेष प्रश्नों का उत्तर लगभग 1000 शब्दों में दीजिए। दाएँ हाशिये में दिए गए अंक अंकों को दर्शाते हैं। अपने उत्तरों को उपयुक्त मानचित्रों और चित्रों द्वारा स्पष्ट कीजिए।}

\bigskip

\noindent\textbf{Question 1.} Write short notes on the following / निम्नलिखित पर संक्षिप्त टिप्पणियाँ लिखें:\pts{4 \times 3.5 = 14}
\begin{parts}
  \item Electro-magnetic Radiation / विद्युत-चुंबकीय विकिरण
  \item Spectral Signature / वर्णक्रमीय हस्ताक्षर
  \item Air borne and Space borne remote sensing / वायु जनित और अंतरिक्ष जनित सुदूर संवेदन
  \item Digital Image Processing / आंकिक प्रतिबिम्ब प्रसंस्करण
\end{parts}

\medskip\rule{\linewidth}{0.2pt}\medskip

\noindent\textbf{Question 2.} Explain the concept and scope of Remote Sensing.\\
\textit{सुदूर संवेदन की अवधारणा और विषय क्षेत्र की व्याख्या कीजिये।}\pts{14}

\medskip\rule{\linewidth}{0.2pt}\medskip

\noindent\textbf{Question 3.} Discuss in detail the spectral bands with suitable examples.\\
\textit{उपयुक्त उदाहरणों सहित वर्णक्रमीय बैंडों पर विस्तृत चर्चा कीजिये।}\pts{14}

\medskip\rule{\linewidth}{0.2pt}\medskip

\noindent\textbf{Question 4.} What do you mean by Remote Sensing Satellite Platforms and sensors? Elaborate with suitable examples.\\
\textit{विभिन्न सुदूर संवेदन उपग्रह प्लेटफॉर्मों और संवेदकों से आप क्या समझते हैं? उपयुक्त उदाहरणों द्वारा स्पष्ट कीजिए।}\pts{14}

\medskip\rule{\linewidth}{0.2pt}\medskip

\noindent\textbf{Question 5.} Discuss the types and characteristics of Aerial photos.\\
\textit{हवाई चित्रों के प्रकार एवं विशेषताओं की विवेचना कीजिए।}\pts{14}

\medskip\rule{\linewidth}{0.2pt}\medskip

\noindent\textbf{Question 6.} Differentiate between Visual and Digital Image Processing techniques.\\
\textit{दृश्य और डिजिटल छवि प्रसंस्करण तकनीकों के बीच अंतर स्पष्ट कीजिए।}\pts{14}

\medskip\rule{\linewidth}{0.2pt}\medskip

\noindent\textbf{Question 7.} Discuss and analyze various Remote Sensing applications in Resource Mapping and Environmental Monitoring.\\
\textit{संसाधन मानचित्रण और पर्यावरण निगरानी में विभिन्न सुदूर संवेदन के अनुप्रयोगों की विवेचना और विश्लेषण कीजिए।}\pts{14}

\vfill
\rule{\linewidth}{0.4pt}
\begin{center}\small
\href{https://bhu-exam.vercel.app/}{Click here to visit our website}\\[4pt]
\href{https://me-aryan.vercel.app/}{Click here to visit our contributor}
\end{center}

\end{document}
